% ==================================================================
%  preamble.tex : bilingual (English | Italiano) geometry booklet
%  Register (agreed with the author): simplified Feynman style,
%  English column for the tutor, Italian column for the student.
% ==================================================================
\documentclass[10pt,a4paper,oneside]{book}

% ---------------------------------------------------------- fonts
\usepackage[T1]{fontenc}
\usepackage[utf8]{inputenc}
\usepackage{amsmath}
\usepackage{libertinus}          % text
\usepackage{libertinust1math}    % matching math
\usepackage[italian,english]{babel}

% ---------------------------------------------------------- layout
\usepackage{geometry}
\geometry{a4paper, left=1.7cm, right=1.7cm, top=2.1cm, bottom=2.0cm,
          headheight=15pt, headsep=10pt, footskip=24pt}
\usepackage{setspace}
\setstretch{1.06}
\setlength{\parindent}{0pt}
\setlength{\parskip}{4pt plus 1pt}
\usepackage{microtype}

% ---------------------------------------------------------- colours
\usepackage[table]{xcolor}
\definecolor{navy}{HTML}{1B2A4A}
\definecolor{crimson}{HTML}{A51C30}
\definecolor{slate}{HTML}{4A4A4A}
\definecolor{forest}{HTML}{23703F}
\definecolor{deeporange}{HTML}{C45A12}
\definecolor{deeppurple}{HTML}{5B2A86}
\definecolor{itgreen}{HTML}{1D6B45}   % accent of the Italian column
\definecolor{enblue}{HTML}{1B2A4A}    % accent of the English column
\definecolor{lightblue}{HTML}{EAF0FA}
\definecolor{lightgray}{HTML}{F2F2F2}
\definecolor{lightgreen}{HTML}{E9F5EC}
\definecolor{lightyellow}{HTML}{FFF8E1}
\definecolor{lightorange}{HTML}{FFF1E3}
\definecolor{lightpurple}{HTML}{F3ECF8}
\definecolor{tutorbg}{HTML}{F4F6F8}
% figure palette (used only inside TikZ)
\definecolor{angA}{HTML}{C0392B}   % red
\definecolor{angB}{HTML}{1F6FB2}   % blue
\definecolor{angC}{HTML}{2E8B57}   % green
\definecolor{angD}{HTML}{D68910}   % amber
\definecolor{shade}{HTML}{DCE6F2}

% ---------------------------------------------------------- packages
\usepackage{tikz}
\usetikzlibrary{arrows.meta, calc, angles, quotes, positioning,
                decorations.pathreplacing, decorations.markings,
                patterns, intersections, shapes.geometric}
\usepackage[most]{tcolorbox}
\usepackage{paracol}
\usepackage{booktabs, tabularx, array, multirow}
\usepackage{graphicx}
\usepackage{caption}
\usepackage{enumitem}
\usepackage{fancyhdr}
\usepackage{titlesec}
\usepackage{ifthen}
\usepackage{needspace}
\usepackage[hidelinks]{hyperref}

\setlist{nosep, leftmargin=1.4em, topsep=2pt}

% ---------------------------------------------------------- paracol
\setlength{\columnsep}{0.75cm}
\setlength{\columnseprule}{0.4pt}
\colseprulecolor{black!20}
\globalcounter*             % figure/equation counters shared by both columns

% language state (global so tcolorbox titles know which column they sit in)
\newif\ifIT
% \begin{bi} English ... \IT Italiano ... \EN English ... \IT ... \end{bi}
\newenvironment{bi}%
  {\begin{paracol}{2}\global\ITfalse\selectlanguage{english}}%
  {\end{paracol}}
\newcommand{\IT}{\switchcolumn\global\ITtrue\selectlanguage{italian}}
\newcommand{\EN}{\switchcolumn*\global\ITfalse\selectlanguage{english}}

% column heads, placed once at the top of each chapter
\newcommand{\colheads}{%
  {\small\sffamily\bfseries\color{enblue}ENGLISH \textperiodcentered\ for the tutor}\par\vspace{-2pt}
  {\color{enblue}\rule{\linewidth}{0.8pt}}%
  \IT
  {\small\sffamily\bfseries\color{itgreen}ITALIANO \textperiodcentered\ per la studentessa}\par\vspace{-2pt}
  {\color{itgreen}\rule{\linewidth}{0.8pt}}%
  \EN}

% ---------------------------------------------------------- headings
\titleformat{\chapter}[display]
  {\normalfont\bfseries\color{navy}}
  {\LARGE\sffamily Chapter\,/\,Capitolo \thechapter}
  {6pt}
  {\Huge}
\titlespacing*{\chapter}{0pt}{-10pt}{14pt}
\titleformat{\section}
  {\normalfont\Large\bfseries\color{navy}}{\thesection}{0.7em}{}
\titleformat{\subsection}
  {\normalfont\large\bfseries\color{slate}}{\thesubsection}{0.7em}{}
\titlespacing*{\section}{0pt}{14pt}{6pt}
\titlespacing*{\subsection}{0pt}{10pt}{4pt}

% bilingual headings:  \bichapter{English}{Italiano}
\newcommand{\bichapter}[2]{%
  \chapter[#1 \textperiodcentered\ #2]{#1\\[2pt]{\color{itgreen}\itshape #2}}}
\newcommand{\bisection}[2]{%
  \Needspace{7\baselineskip}%
  \section[#1 \textperiodcentered\ #2]{#1 {\color{itgreen}\itshape\,\textperiodcentered\ #2}}}
\newcommand{\bisubsection}[2]{%
  \Needspace{6\baselineskip}%
  \subsection[#1 \textperiodcentered\ #2]{#1 {\color{itgreen}\itshape\,\textperiodcentered\ #2}}}

% ---------------------------------------------------------- header/footer
\pagestyle{fancy}
\fancyhf{}
\fancyhead[L]{\small\color{slate}\nouppercase{\leftmark}}
\fancyhead[R]{\small\color{slate}Angles \& Congruence \textperiodcentered\ \textit{Angoli e congruenza}}
\fancyfoot[C]{\small\thepage}
\renewcommand{\headrulewidth}{0.4pt}
\renewcommand{\chaptermark}[1]{\markboth{\thechapter.\ #1}{}}
\fancypagestyle{plain}{\fancyhf{}\fancyfoot[C]{\small\thepage}\renewcommand{\headrulewidth}{0pt}}

% ---------------------------------------------------------- captions
\captionsetup{font=small, labelfont={bf,color=navy}, labelsep=period,
              justification=centering, skip=4pt}
\captionsetup[figure]{name={Figure\,/\,Figura}}
\captionsetup[table]{name={Table\,/\,Tabella}}
% \bicap{English caption}{Didascalia italiana}
\newcommand{\bicap}[2]{\caption{#1\\{\color{itgreen}\itshape #2}}}

% ---------------------------------------------------------- math shorthands
\newcommand{\ang}[3]{#1\widehat{#2}#3}      % angle ABC  ->  A \hat B C
\newcommand{\av}[1]{\widehat{#1}}           % angle named by its vertex
\newcommand{\seg}[1]{\overline{#1}}         % segment
\newcommand{\tri}[1]{\triangle #1}          % triangle
\newcommand{\gradi}{^\circ}

% ---------------------------------------------------------- boxes
% Every box is FULL WIDTH and split in two like the page:
%   \begin{trap}{English subtitle}{Sottotitolo italiano}
%     English text
%   \tcblower
%     Testo italiano
%   \end{trap}
% Either subtitle may be empty.
\newlength{\halfcol}
\newcommand{\bxhalf}[3]{\if\relax\detokenize{#2}\relax#1\else#1: #2\fi}
% title bar split at the same place as the box body
\newcommand{\bxt}[4]{%
  \setlength{\halfcol}{\dimexpr(\linewidth-\columnsep)/2\relax}%
  \parbox[t]{\halfcol}{\raggedright\bxhalf{#1}{#3}{}}\hfill
  \parbox[t]{\halfcol}{\raggedright\bxhalf{#2}{#4}{}}}
\tcbset{bookbase/.style={enhanced, sharp corners=south, arc=2.5pt,
        boxrule=0.8pt, left=5pt, right=5pt, top=3pt, bottom=4pt,
        fonttitle=\sffamily\bfseries\small, coltitle=white,
        before skip=7pt plus 2pt, after skip=7pt plus 2pt, toptitle=1pt, bottomtitle=1pt,
        sidebyside, sidebyside align=top seam, sidebyside gap=\columnsep,
        lefthand ratio=0.5,
        before upper={\global\ITfalse\selectlanguage{english}},
        before lower={\global\ITtrue\selectlanguage{italian}},
        after={\global\ITfalse\selectlanguage{english}}}}

\newtcolorbox{trap}[2]{bookbase, colback=lightgray, colframe=slate,
  segmentation style={slate!40, solid},
  title={\bxt{The trap}{La trappola}{#1}{#2}}}
\newtcolorbox{rebuild}[2]{bookbase, colback=lightblue, colframe=navy,
  segmentation style={navy!30, solid},
  title={\bxt{Step by step}{Passo dopo passo}{#1}{#2}}}
\newtcolorbox{anchor}[2]{bookbase, colback=lightgreen, colframe=forest,
  segmentation style={forest!30, solid},
  title={\bxt{Picture it}{Immagina}{#1}{#2}}}
\newtcolorbox{think}[2]{bookbase, colback=lightorange, colframe=deeporange,
  segmentation style={deeporange!30, solid},
  title={\bxt{Think first}{Pensaci prima}{#1}{#2}}}
\newtcolorbox{example}[2]{bookbase, colback=lightpurple, colframe=deeppurple,
  segmentation style={deeppurple!30, solid},
  title={\bxt{Worked example}{Esempio svolto}{#1}{#2}}}
\newtcolorbox{keyidea}[2]{bookbase, colback=white, colframe=crimson, boxrule=1.3pt,
  segmentation style={crimson!30, solid}, fontupper=\itshape, fontlower=\itshape,
  title={\bxt{The click}{Il clic}{#1}{#2}}}
\newtcolorbox{defn}[2]{bookbase, colback=white, colframe=navy, boxrule=0.8pt,
  segmentation style={navy!30, solid},
  title={\bxt{Definition}{Definizione}{#1}{#2}}}
\newtcolorbox{history}[2]{bookbase, colback=lightyellow, colframe=crimson!80!black,
  segmentation style={crimson!30, solid},
  title={\bxt{A bit of history}{Un po' di storia}{#1}{#2}}}
\newtcolorbox{proofbox}[2]{bookbase, colback=white, colframe=navy!70,
  segmentation style={navy!30, solid},
  title={\bxt{Proof}{Dimostrazione}{#1}{#2}}}
\newtcolorbox{exercise}[2]{bookbase, colback=white, colframe=forest!80!black,
  segmentation style={forest!30, solid},
  title={\bxt{Exercise}{Esercizio}{#1}{#2}}}
\newtcolorbox{solution}[2]{bookbase, colback=white, colframe=slate!80,
  segmentation style={slate!30, solid},
  title={\bxt{Solution}{Soluzione}{#1}{#2}}}

% full width note for the tutor only (outside the parallel columns)
\newtcolorbox{tutor}[1][]{enhanced, colback=tutorbg, colframe=navy!40, boxrule=0.5pt,
  borderline west={2.5pt}{0pt}{navy}, sharp corners, left=7pt, right=6pt, top=3pt, bottom=3pt,
  fontupper=\small, before skip=8pt, after skip=8pt,
  title={For the tutor\if\relax\detokenize{#1}\relax\else: #1\fi},
  fonttitle=\sffamily\bfseries\small, coltitle=navy, colbacktitle=tutorbg,
  attach title to upper, after title={.\ }}

% "answer" line inside think boxes
\newcommand{\answer}{\par\vspace{2pt}{\color{deeporange}\rule[0.5ex]{\linewidth}{0.4pt}}\par
  {\sffamily\bfseries\small\color{deeporange}\ifIT Risposta\else Answer\fi.}\ }

% statement / reason proof lines
\newcommand{\step}[2]{\item #1\hfill{\small\color{slate}\itshape(#2)}}
\newenvironment{steps}{\begin{enumerate}[label=\arabic*., leftmargin=1.5em, itemsep=1pt]}{\end{enumerate}}

% ---------------------------------------------------------- TikZ styles
\tikzset{
  pt/.style={circle, fill=black, inner sep=1.1pt},
  lbl/.style={font=\small},
  side/.style={line width=0.9pt, line cap=round, line join=round},
  thinline/.style={line width=0.5pt},
  helper/.style={line width=0.5pt, dashed, black!55},
  ray/.style={line width=0.9pt, -{Stealth[length=2.2mm]}},
  rightangle/.style={draw, line width=0.5pt},
  % tick marks for congruent segments
  tick1/.style={postaction={decorate, decoration={markings,
      mark=at position 0.5 with {\draw[line width=0.7pt] (0,-3pt)--(0,3pt);}}}},
  tick2/.style={postaction={decorate, decoration={markings,
      mark=at position 0.5 with {\draw[line width=0.7pt] (-1.2pt,-3pt)--(-1.2pt,3pt) (1.2pt,-3pt)--(1.2pt,3pt);}}}},
  tick3/.style={postaction={decorate, decoration={markings,
      mark=at position 0.5 with {\draw[line width=0.7pt] (-2.2pt,-3pt)--(-2.2pt,3pt) (0,-3pt)--(0,3pt) (2.2pt,-3pt)--(2.2pt,3pt);}}}},
}
% right angle mark at vertex #2 between points #1 and #3:  \rang{A}{H}{B}
\newcommand{\rang}[3]{%
  \pic[draw, line width=0.5pt, angle radius=2.2mm] {right angle=#1--#2--#3};}
